\subsection{\texorpdfstring{\(W_{a}\) 的正规化子}{W\_a 的正规化子}}

本号中，我们假设 V 上所选的标量积不仅在 W 下不变，而且在整个群 A(R) 下不变。我们认同 A(R) 与 A(R \(^{\prime}\)）。

令 G 为 \(W_{a}\) 在仿射欧几里得空间 E 的位移群中的正规化子。若 g 是 E 的一个位移，s 是关于超平面 L 的正交反射，则位移 \(gsg^{-1}\) 是关于超平面 \(g(L)\) 的正交反射。因此，G 正是置换超平面 \(L_{\alpha,k}\)（其中 \(\alpha \in R\) 且 \(k \in Z\)）的 E 的位移所成的集合。

现在，E 的自同构群是 \(V^{*}\) 的正交群 U 与平移群 T 的半直积。若 \(u \in U\) 且 \(v \in V^{*}\)，则超平面 \(L_{\alpha,k}\) 在 \(g = u \circ t(v)\) 下变为方程为

\[
\langle^ {t} u ^ {- 1} (\alpha), x \rangle = k + \langle \alpha , v \rangle .
\]

因此，\(g \in \mathbf{G}\) 当且仅当，一方面 \({}^t u\) 置换这些根，亦即属于 \(\mathrm{A(R)}\)；另一方面有 \(\langle \alpha, v \rangle \in \mathbf{Z}\)（对所有 \(\alpha \in \mathbb{R}\)），即 \(v \in \mathrm{P(R^{\prime})}\)。换言之，群 \(G\) 是 \(A(R)\) 与 \(P\) 的半直积。由于 \(Q \subseteq P\) 且 \(W \subseteq A(R)\)，商群 \(G / W_a\) 是 \(A(R)/W\) 与 \(P(R^{\prime})/Q(R^{\prime})\) 的半直积；容易验证，\(A(R)/W\) 在 \(P(R^{\prime})/Q(R^{\prime})\) 上的相应作用就是典范作用（第 1 节，第 9 号）。

记 \(W_{a}^{\prime}\) 为由 W 与 P 的半直积构成的 G 的子群。这是 G 的正规子群，且 \(G/W_{a}^{\prime}\) 典范同构于 \(A(R)/W\)；此外，从 \(P(R^{\prime})\) 到 \(W_{a}^{\prime}/W_{a}\) 的典范映射通过取商给出 \(P(R^{\prime})/Q(R^{\prime})\) 到 \(W_{a}^{\prime}/W_{a}\) 的一个同构。

现在令 C 为 E 的一个小室，令 \(G_{C}\) 为由元素 \(g \in G\)（满足 \(g(\mathrm{C}) = \mathrm{C}\)）所成的子群。由于 \(W_{a}\) 在小室上单纯传递地作用，群 G 是 \(G_{C}\) 与 \(W_{a}\) 的半直积。相应的从 \(G/W_{a}\) 到 \(G_{C}\) 的同构尤其给出从 \(\mathrm{P}(\mathrm{R}^{\prime})/\mathrm{Q}(\mathrm{R}^{\prime})\) 到群 \(\Gamma_{C} = G_{C} \cap W_{a}^{\prime}\) 的一个典范同构。

假设 R 不可约，并保留第 2 号命题 5 的记号。置 \(R_{0} = R\)，并令 \(R_{i} (i \in I)\) 为由 \(a_{j}\)（\(j \neq i\)）生成的根系。对于 i = 0（分别对于 \(i \in I\)），令 \(w_{i}\) 为 \(W(R_{i})\) 中的唯一元素（将其视为 W 的子群），它把 \(R_{i}\) 相对于基 \((\alpha_{j})_{j \neq i}\) 的正根变为负根（第 1 节，第 6 号，命题 17 的推论 3）。

命题 6. 对所有 \(i \in \mathbf{J}\)，元素 \(\gamma_i = t(\overline{\omega}_i)w_i w_0\) 属于 \(\Gamma_{C}\)，且映射 \(i \mapsto \gamma_i\) 是从 \(\mathbf{J}\) 到 \(\Gamma_{C} - \{1\}\) 的双射。

首先注意，根 \(w_{i}(\tilde{\alpha})\) 具有形式

\[
n _ {i} \alpha_ {i} + \sum_ {j \neq i} b _ {i j} \alpha_ {j},
\]

因而是正根。

我们证明，若 \(i \in \mathrm{J}\)，则 \(\gamma_i \in \Gamma_{C}\)。事实上，设 \(a \in \mathrm{C}\) 且 \(b = \gamma_i(a)\)。对于 \(1 \leqslant j \leqslant l\) 且 \(j \neq i\)，

\[
\begin{array}{r l} \langle b, \alpha_ {j} \rangle & = \langle \overline {{\omega}} _ {i} + w _ {i} w _ {0} (a), \alpha_ {j} \rangle \\ & = \langle w _ {0} (a), w _ {i} (\alpha_ {j}) \rangle > 0 \end{array}\tag{2}
\]

因为 \(w_0(a) \in -\mathbf{C}'\)，且 \(w_i(\alpha_j)\) 为负根。另一方面，

\[
\langle b, \alpha_ {i} \rangle = 1 + \left\langle w _ {0} (a), w _ {i} (\alpha_ {i}) \right\rangle \geqslant 1 + \left\langle w _ {0} (a), \tilde {\alpha} \right\rangle > 0\tag{3}
\]

因为 \(w_0(a) \in -\mathbf{C}'\)，\(\tilde{\alpha} - w_i(\alpha_i)\) 在 \(-\mathbf{C}'\) 上取负值，且 \(\langle w_0(a), \tilde{\alpha} \rangle > -1\)。最后，

\[
\langle b, \tilde {\alpha} \rangle = n _ {i} + \left\langle w _ {0} (a), w _ {i} (\tilde {\alpha}) \right\rangle = 1 + \left\langle w _ {0} (a), w _ {i} (\tilde {\alpha}) \right\rangle <   1\tag{4}
\]

因为 \(w_0(a) \in -\mathbf{C}'\)，且 \(w_i(\tilde{\alpha})\) 是正根。于是关系 (2)、(3) 和 (4) 蕴含 \(b \in \mathbf{C}\)，从而 \(\gamma_i \in \Gamma_{C}\)。显然，映射 \(i \mapsto \gamma_i\) 是单射，因为 \(\gamma_i(0) = \overline{\omega}_i\)。最后，设 \(\gamma \in \Gamma_{C}\) 且 \(\gamma \neq 1\)，并令 \(\gamma = tw\)，其中 \(t \in \mathbf{P}\) 且 \(w \in W\)。于是 \(t \neq 1\)，因为 \(\Gamma_{C} \cap W = \{1\}\)。另一方面，\(t(0) = \gamma(0) \in \overline{\mathbf{C}} \cap P(\mathbf{R}^{\prime})\)，命题 5 蕴含存在 \(i \in J\) 使得 \(t(0) = \overline{\omega}_i\)。于是 \(\gamma_i^{-1}\gamma(0) = 0\)，故 \(\gamma = \gamma_i\)，因为 \(\Gamma_{C} \cap W = \{1\}\)。证明完毕。

推论。 \((\bar{\omega}_i)_{i\in J}\) 在 \(\mathrm{P}(\mathbb{R}^{\prime})\) 中构成 \(\mathrm{P}(\mathbb{R}^{\prime}) / \mathrm{Q}(\mathbb{R}^{\prime})\) 的非零元素的一组代表元。

事实上，若将 \(\Gamma_{C}\) 与 \(\mathrm{P}(\mathbb{R}^{\prime}) / \mathrm{Q}(\mathbb{R}^{\prime})\) 认同，则元素 \(\gamma_{i}\) 对应于 \(\overline{\omega}_{i}\) 模 \(\mathrm{Q}(\mathbb{R}^{\prime})\) 的陪集。

备注。1）映射 \(\gamma \mapsto \gamma(0)\) 是从 \(\Gamma_{C}\) 到 \(\overline{\mathrm{C}} \cap \mathrm{P}(\mathrm{R}^{\prime})\) 的双射。

\begin{enumerate}
\def\labelenumi{\arabic{enumi})}
\setcounter{enumi}{1}
\tightlist
\item
  若 E 只赋予其仿射结构，则群 G 也是 \(W_{a}\) 在 E 的自同构群中的正规化子（参见习题 3）。
\end{enumerate}

